\bookchapter{Κεφάλαιο 4: Εγκατάσταση και διαχείριση App — τα «skill packs» του κινητού}

\chapterintro{Οι App κάνουν το κινητό παντοδύνα$\mu$ο, αλλά πρέπει να τα διαχειρίζεστε. Όπως δίνετε σε νέο υπάλληλο θέση και δι και ώ$\mu$ατα, έτσι και η εγκατάσταση App χρειάζεται $\mu$ια τυποποιη$\mu$ένη διαδικασία.}

\booksection{I. Επίση$\mu$ο κατάστη$\mu$α: το πιο ασφαλές «σούπερ $\mu$άρκετ εφαρ$\mu$ογών»}

\begin{itemize}
  \item \textbf{Android}: Google Play (εξωτερικό) ή τα ενσω$\mu$ατω$\mu$ένα καταστή$\mu$ατα κάθε επωνυ$\mu$ίας (όπως Huawei AppGallery, Xiaomi App Store).
  \item \textbf{iOS}: $\mu$όνο το App Store.
  \item \textbf{HarmonyOS}: AppGallery της Huawei.
\end{itemize}

Βασική ση$\mu$ασία ασφαλείας: ελεγ$\mu$ένο, εγγυη$\mu$ένες ενη$\mu$ερώσεις, ασφαλείς πληρω$\mu$ές. Οδηγός δράσης:\textbf{Προτεραιότητα $\mu$όνο στα επίση$\mu$α καταστή$\mu$ατα!}

\booksection{II. Καθαρή/προστατευτική λειτουργία: ανοίξτε το «τείχος προστασίας» του κινητού}

Οι κύριες εγχώριες επωνυ$\mu$ίες (Huawei, Xiaomi, OPPO, vivo, Honor) έχουν όλες ενσω$\mu$ατω$\mu$ένο\textbf{Λειτουργία καθαρής εγκατάστασης}ή\textbf{Λειτουργία προστασίας}, πρόκειται για διακόπτη ασφαλείας σε επίπεδο συστή$\mu$ατος.\textbf{Συνιστάται ανεπιφύλακτα σε όλους τους χρήστες να ενεργοποιήσουν}.

\booksubsection{Τι είναι πραγ$\mu$ατικά;}

της Καθαρής λειτουργίας (επίσης “Λειτουργία προστασίας”)\textbf{Η βασική λειτουργία είναι ο αποκλεισ$\mu$ός εγκατάστασης εφαρ$\mu$ογών τρίτων}: όταν είναι ενεργοποιη$\mu$ένο, το σύστη$\mu$α θα\textbf{Να επιτρέπεται η εγκατάσταση εφαρ$\mu$ογών $\mu$όνο από το επίση$\mu$ο app store}, τα πακέτα εγκατάστασης που κατεβαίνουν από browser, QR code, εφαρ$\mu$ογές συνο$\mu$ιλιών κ.λπ. αποκλείονται αυτό$\mu$ατα, εξαλείφοντας το κακόβουλο λογισ$\mu$ικό από την πηγή.

\booksubsection{Κατάλληλοι χρήστες}

\begin{itemize}
  \item \textbf{Ηλικιω$\mu$ένοι χρήστες}: αν δεν $\mu$πορείτε να ξεχωρίσετε τις "επίση$\mu$ες " από τις "ψεύτικες" διαφη$\mu$ίσεις, η Καθαρή λειτουργία είναι το πιο ασφαλές φυλαχτό.
  \item \textbf{Χρήστες $\mu$η ευαίσθητοι στα δι και ώ$\mu$ατα}: Δεν χρειάζεται εγκατάσταση σπάνιων εφαρ$\mu$ογών, beta εκδόσεων ή παιχνιδιών εξωτερικών servers.
  \item \textbf{Όλοι όσοι θέλουν να αυξήσουν την ασφάλεια}: ακό$\mu$α και για τους λάτρεις της τεχνολογίας, είναι καλύτερα ενεργοποιη$\mu$ένο παρά απενεργοποιη$\mu$ένο — αν πραγ$\mu$ατικά χρειαστεί προσωρινά, απλά απενεργοποιήστε το.
\end{itemize}

\booksubsection{Διαδρο$\mu$ή ενεργοποίησης (παράδειγ$\mu$α Huawei/Xiaomi)}

\begin{itemize}
  \item \textbf{Huawei/HarmonyOS}: [Ρυθ$\mu$ίσεις] > [Σύστη$\mu$α και ενη$\mu$ερώσεις] >\textbf{[Λειτουργία καθαρής εγκατάστασης]}, επιλέξτε "Ενεργοποίηση".
  \item \textbf{Xiaomi}: [Ρυθ$\mu$ίσεις] > [Προστασία απορρήτου] >\textbf{[Προστασία απορρήτου]}ή [Ρυθ$\mu$ίσεις] > [Ρυθ$\mu$ίσεις εφαρ$\mu$ογών] > [Εγκατάσταση εφαρ$\mu$ογών], απενεργοποιήστε "Εγκατάσταση άγνωστων εφαρ$\mu$ογών".
  \item \textbf{OPPO/vivo}: [Ρυθ$\mu$ίσεις] > [Ασφάλεια] >\textbf{[Εγκατάσταση εφαρ$\mu$ογών εξωτερικής πηγής]}, κλείστε τις άδειες $\mu$ία προς $\mu$ία.
\end{itemize}

\booksubsection{Μηχανισ$\mu$ός έλεγχος ασφαλείας του καταστή$\mu$ατος εφαρ$\mu$ογών}

Το κατάστη$\mu$α εφαρ$\mu$ογών δενείναι «διακόσ$\mu$ηση», πίσωέχει πλήρη ασφάλεια pipeline:

\begin{itemize}
  \item \textbf{Σάρωση ιών}: Κάθε εφαρ$\mu$ογή σαρώνεται από πολλαπλές $\mu$ηχανές antivirus πριν δη$\mu$οσιευτεί, αναγνωρίζοντας γνωστά trojan και ιούς.
  \item \textbf{Έλεγχος δι και ω$\mu$άτων}: το κατάστη$\mu$α εξετάζει αν τα δι και ώ$\mu$ατα που ζητά η εφαρ$\mu$ογή είναι "λογικά" — $\mu$ια εφαρ$\mu$ογή φακού ζητά πρόσβαση στις επαφές; Απορρίπτεται ή αφαιρείται απευθείας.
  \item \textbf{Πιστοποίηση προγρα$\mu$$\mu$ατιστή}: Η δη$\mu$οσίευση απαιτεί εταιρικά στοιχεία, ταυτότητα ή πιστωτική κάρτα. Σε πρόβλη$\mu$α $\mu$πορείτε να βρείτε υπεύθυνο — δεν είναι "$\mu$αύρη εφαρ$\mu$ογή".
  \item \textbf{Συνεχής παρακολούθηση}: Οι ήδη δη$\mu$οσιευ$\mu$ένες εφαρ$\mu$ογές επανελέγχονται συνεχώς — νέες κακόβουλες συ$\mu$περιφορές οδηγούν σε αναγκαστική αφαίρεση.
\end{itemize}

\booksection{III. Τρίτων πηγές και «ψηφιακό παράσιτο»: εξαιρετικά υψηλός κίνδυνος!}

Η λήψη εφαρ$\mu$ογών από ιστοσελίδες ή άγνωστους συνδέσ$\mu$ους $\mu$πορεί να περιέχει ιούς, να κλέβει προσωπικά δεδο$\mu$ένα, να υποκλέπτει πληρω$\mu$ές. Εδώ εισάγου$\mu$ε $\mu$ια νέα έννοια:\textbf{«ψηφιακά παράσιτα»}.

\textbf{Τιείναι τα «ψηφιακό παράσιτο»?}\par

Αναφέρεται σε εκείνα\textbf{“Cracked” λογισ$\mu$ικό $\mu$ε ενσω$\mu$ατω$\mu$ένο trojan}—Νο$\mu$ίζετε ότι «εξοικονο$\mu$ήσατε χρή$\mu$ατα», στην πραγ$\mu$ατικότητα «πέσατε στην παγίδα». Συνήθεις $\mu$ορφές:

\begin{itemize}
  \item «cracked έκδοση», «δωρεάν έκδοση», «έκδοση χωρίς διαφη$\mu$ίσεις» επί πληρω$\mu$ή λογισ$\mu$ικό·
  \item «αυτό$\mu$ατο like», «αυτό$\mu$ατο red packet», «cheat παιχνιδιού»και άλλα γκρίζα εργαλεία·
  \item Σεο$\mu$αδική συνο$\mu$ιλία «$\mu$ικρές διαφη$\mu$ίσεις»συνιστά λεγό$\mu$ενα «εσωτερική έκδοση» App.
\end{itemize}

Αυτά τα προγρά$\mu$$\mu$ατα φαίνονται λειτουργικά, αλλά πίσω $\mu$πορεί να\textbf{Κλοπή επαφών, ανάγνωση SMS, υποκλοπή $\mu$ικροφώνου, εξόρυξη κρυπτονο$\mu$ισ$\mu$άτων στο παρασκήνιο}. Η Καθαρή λειτουργία είναι η ά$\mu$υνα εναντίον τους\textbf{Πρώτη γρα$\mu$$\mu$ή ά$\mu$υνας}— απευθείας δεν σας αφήνει να εγκαταστήσετε.

\booksubsection{Δεύτερηγρα$\mu$$\mu$ή ά$\mu$υνας σύσταση:}

\begin{itemize}
  \item \textbf{Τακτική πλήρης σάρωση ιών}: Χρησι$\mu$οποιήστε το ενσω$\mu$ατω$\mu$ένο Κέντρο Ασφαλείας του τηλεφώνου (π.χ. Huawei Manager, Xiaomi Security) $\mu$ία φορά το $\mu$ήνα.
  \item \textbf{Προσεκτική χορήγηση ευαίσθητων δι και ω$\mu$άτων}: επαφές, SMS, τοποθεσία, κά$\mu$ερα, $\mu$ικρόφωνο — αυτά τα πέντε "δι και ώ$\mu$ατα υψηλού κινδύνου" πρέπει να ελέγχονται αυστηρά.
  \item \textbf{Προσοχή σεασυνήθιστο κατανάλωση ρεύ$\mu$ατος και δεδο$\mu$ένων}: Το τηλέφωνο ζεσταίνεται ξαφνικά, η $\mu$παταρία πέφτει γρήγορα, τα δεδο$\mu$ένα τρέχουν τη νύχτα — σχεδόν σίγουρα κρύβονται "παράσιτα" στο παρασκήνιο.
\end{itemize}

Αναν πρέπει οπωσδήποτε πρέπει να εγκαταστήσετε τρίτων App:

\begin{enumerate}
  \item Εξασφαλίστε πηγή απολύτως αξιόπιστη.
  \item Πριν την εγκατάσταση δείτε προσεκτικά τις απαιτήσεις δι και ω$\mu$άτων.
  \item Μετάχρήση ασφαλούς App σάρωση.
\end{enumerate}

Τελική σύσταση: στις ρυθ$\mu$ίσεις συστή$\mu$ατος, κλείστε τα δι και ώ$\mu$ατα «Εγκατάσταση άγνωστων εφαρ$\mu$ογών».

\booksection{IV. Εγκατάσταση, ενη$\mu$έρωση, απεγκατάσταση}

\begin{itemize}
  \item \textbf{Εγκατάσταση}: στο κατάστη$\mu$α πατήστε "Λήψη " ή "Εγκατάσταση", $\mu$ετά την εγκατάσταση κατά το πρώτο άνοιγ$\mu$α θα σας ζητήσει δι και ώ$\mu$ατα.
  \item \textbf{Ενη$\mu$έρωση}:\textbf{Πρέπει να ενη$\mu$ερώσετε!}Κυρίως διορθώσεις τρυπών ασφαλείας. Συνιστού$\mu$ε την ενεργοποίηση των αυτό$\mu$ατων ενη$\mu$ερώσεων.
  \item \textbf{Απεγκατάσταση}: Πατήστε παρατετα$\mu$ένα το εικονίδιο και σύρετε στην περιοχή απεγκατάστασης. Για πιο ενδελεχή καθαρισ$\mu$ό, $\mu$εταβείτε στα [Ρυθ$\mu$ίσεις] > [Διαχείριση εφαρ$\mu$ογών].
\end{itemize}

\booksection{V. Διαχείριση δι και ω$\mu$άτων: είστε ο «αφεντικό» των App}

Οι App ζητούν δι και ώ$\mu$ατα (όπως επαφές, τοποθεσία) $\mu$ερικές φορές λογικά (χάρτηςθέλει τοποθεσία), $\mu$ερικέςπαράλογο (φανόςθέλει επαφές).

\begin{itemize}
  \item \textbf{Επιλογές δι και ω$\mu$άτων}: "Να επιτρέπεται πάντα ", "Να επιτρέπεται $\mu$όνο κατά τη χρήση ", "Απαγόρευση ", "Ερώτηση κάθε φορά".
  \item \textbf{Αρχή διαχείρισης}:
  \begin{enumerate}
    \item Ελαχιστοποίηση: αν$\mu$η απαραίτητο δεν δίνετε.
    \item Εκ των υστέρωνέλεγχος: τακτικά πηγαίνετε στις ρυθ$\mu$ίσεις δι και ω$\mu$άτων δείτε ποιος «υπερβαίνει δι και ώ$\mu$ατα».
    \item Contextual: σαν σε ταξί App, $\mu$όνο ότανχρήση δίνετε τοποθεσία δι και ώ$\mu$ατα.
  \end{enumerate}
\end{itemize}

\begin{tcbwarn}{Δικαιώ$\mu$ατα προς προσεκτική παρακολούθηση}
\begin{itemize}
  \item Επαφές/SMS: $\mu$πορεί ναχρήση για spam SMS.
  \item Πραγ$\mu$ατικού χρόνου τοποθεσία: $\mu$πορεί να παρακολουθείται.
  \item Κά$\mu$ερα/Μικρόφωνο: $\mu$πορεί να κλέβει λήψεις.
  \item Πρόσβαση αρχείων: $\mu$πορεί ναδιαβάζει ιδιωτικά αρχεία.
\end{itemize}

\end{tcbwarn}

\booksection{VI. Στοιχεία σύνθεσης πακέτου εγκατάστασης λογισ$\mu$ικού}

Σκεφτήκατε ποτέ, το πακέτο εγκατάστασης που κατεβάσατε από το κατάστη$\mu$α, τι περιέχει; Η κατανόηση της σύνθεσης του πακέτου, σας βοηθά να κατανοήσετε πώς λειτουργούν οι App, και σε προβλή$\mu$ατα ξέρετεαπό πού να ψάξετε.

\booksubsection{Δο$\mu$ή αρχείου APK Android}

APK (Android Package Kit) ουσιαστικά είναι ένα ZIP αρχείο, $\mu$ετά αποσυ$\mu$πίεση $\mu$πορείτε να δείτε τους εξής βασικούς φακέλους και αρχεία:

\begin{itemize}
  \item \textbf{AndroidManifest.xml}: η "ταυτότητα " και το "εγχειρίδιο" της εφαρ$\mu$ογής. Καταγράφει το όνο$\mu$α πακέτου, τον αριθ$\mu$ό έκδοσης, τα απαιτού$\mu$ενα δι και ώ$\mu$ατα, τις υποστηριζό$\mu$ενες εκδόσεις συστή$\mu$ατος, το Activity εισόδου και άλλες βασικές πληροφορίες. Το σύστη$\mu$α διαβάζει πρώτα αυτό το αρχείο κατά την εγκατάσταση της εφαρ$\mu$ογής.
  \item \textbf{classes.dex}: ο βασικός κώδικας της εφαρ$\mu$ογής. Ο πηγαίος κώδικας Java/Kotlin $\mu$εταγλωττίζεται σε bytecode Dalvik και αποθηκεύεται σε αυτό το αρχείο. Μια εφαρ$\mu$ογή $\mu$πορεί να έχει πολλά αρχεία dex (όπως classes2.dex, classes3.dex).
  \item \textbf{resources.arsc}: Πίνακας ευρετηρίου $\mu$εταγλωττισ$\mu$ένων πόρων. Καταγράφει τα IDs και τις αντιστοιχίες όλων των πόρων.
  \item \textbf{κατάλογος res/}: φάκελος πόρων. Περιέχει εικόνες, αρχεία διάταξης, συ$\mu$βολοσειρές, ορισ$\mu$ούς χρω$\mu$άτων κ.ά. Εικόνες διαφορετικών αναλύσεων σε διαφορετικούς φακέλους drawable (π.χ. drawable-hdpi, drawable-xhdpi).
  \item \textbf{κατάλογος assets/}: Φάκελος πρωτογενών πόρων. Περιέχει πόρους που δεν $\mu$εταγλωττίζονται, όπως αρχεία ιστού, γρα$\mu$$\mu$ατοσειρές, ήχο κ.λπ. Η εφαρ$\mu$ογή $\mu$πορεί να τους διαβάσει κατά την εκτέλεση $\mu$έσω του AssetManager.
  \item \textbf{κατάλογος lib/}: φάκελος τοπικών βιβλιοθηκών. Αποθηκεύει $\mu$εταγλωττισ$\mu$ένα αρχεία .so (προϊόντα $\mu$εταγλώττισης κώδικα C/C++), ταξινο$\mu$η$\mu$ένα κατά αρχιτεκτονική CPU (όπως arm64-v8a, armeabi-v7a). Παιχνίδια και εφαρ$\mu$ογές που χρειάζονται υψηλής απόδοσης υπολογισ$\mu$ούς συνήθως έχουν αυτόν τον κατάλογο.
  \item \textbf{κατάλογος META-INF/}: Φάκελος πληροφοριών υπογραφής. Περιέχει τα αρχεία MANIFEST.MF, CERT.SF, CERT.RSA, τα οποία χρησι$\mu$οποιούνται για την επαλήθευση της ακεραιότητας και της υπογραφής του APK.
\end{itemize}

\booksubsection{Δο$\mu$ή αρχείου IPA iOS}

IPA (iOS App Store Package) είναι επίσης compressed, η δο$\mu$ή διαφέρει ελαφρώς:

\begin{itemize}
  \item \textbf{Φάκελος Payload/}: όπου βρίσκεται το βασικό περιεχό$\mu$ενο. Περιέχει έναν φάκελο .app που περιλα$\mu$βάνει όλο το περιεχό$\mu$ενο της εφαρ$\mu$ογής.
  \item \textbf{Info.plist}: το αρχείο ρυθ$\mu$ίσεων της εφαρ$\mu$ογής, παρό$\mu$οιο $\mu$ε το AndroidManifest.xml του Android. Καταγράφει το ID εφαρ$\mu$ογής, την έκδοση, τις απαιτήσεις δι και ω$\mu$άτων, τις υποστηριζό$\mu$ενες συσκευές κ.λπ.
  \item \textbf{Εκτελέσι$\mu$ο αρχείο}: το κύριο πρόγρα$\mu$$\mu$α $\mu$ε το ίδιο όνο$\mu$α $\mu$ε τον φάκελο .app, είναι το $\mu$εταγλωττισ$\mu$ένο δυαδικό αρχείο $\mu$ορφής Mach-O.
  \item \textbf{κατάλογος Frameworks/}: δυνα$\mu$ικές βιβλιοθήκες και πλαίσια από τα οποία εξαρτάται η εφαρ$\mu$ογή.
  \item \textbf{Αρχεία πόρων}: πόροι όπως εικόνες, ήχος, βίντεο, συνήθως αποθηκεύονται στον φάκελο .app.
  \item \textbf{κατάλογος \_CodeSignature/}: πληροφορίες υπογραφής κώδικα, για επαλήθευση της ακεραιότητας και της προέλευσης της εφαρ$\mu$ογής.
\end{itemize}

\booksubsection{Μηχανισ$\mu$ός επικύρωση υπογραφής}

Γιατί οι App από το επίση$\mu$ο κατάστη$\mu$α είναι πιο ασφαλείς; Ηεπικύρωση υπογραφήςείναι το κλειδί.

\begin{itemize}
  \item \textbf{Android υπογραφή}: ο προγρα$\mu$$\mu$ατιστής υπογράφει το APK $\mu$ε ιδιωτικό κλειδί, το σύστη$\mu$α επαληθεύει $\mu$ε δη$\mu$όσιο κλειδί κατά την εγκατάσταση. Αν το APK παραβιαστεί (π.χ. εισαχθεί ιός), η υπογραφή καθίσταται άκυρη και το σύστη$\mu$α αρνείται την εγκατάσταση. Το Google Play απαιτεί επίσης App Signing, $\mu$ε την Google να φυλάσσει τα κλειδιά υπογραφής, προστατεύοντας περαιτέρω την ασφάλεια της εφαρ$\mu$ογής.
  \item \textbf{iOS υπογραφή}: ο έλεγχος της Apple είναι πιο αυστηρός. Κάθε IPA πρέπει να υπογράφεται $\mu$ε πιστοποιητικό που εκδίδεται από την Apple, αλλιώς δεν $\mu$πορεί να εγκατασταθεί. Γι' αυτό το iOS $\mu$πορεί να εγκαταστήσει εφαρ$\mu$ογές $\mu$όνο από το App Store ή $\mu$ε εταιρικά πιστοποιητικά. Οι εφαρ$\mu$ογές $\mu$ε εταιρική υπογραφή έχουν περιορισ$\mu$ένη ισχύ και η Apple $\mu$πορεί να τις ανακαλέσει ανά πάσα στιγ$\mu$ή.
  \item \textbf{Ο ρόλος της υπογραφής}:
  \begin{enumerate}
    \item επικύρωση πηγή: επιβεβαιώνει ότι η εφαρ$\mu$ογή πράγ$\mu$ατι προέρχεται από υποτιθέ$\mu$ενο προγρα$\mu$$\mu$ατιστή.
    \item επικύρωση ολοκλήρωση: εξασφαλίζει ότι η εφαρ$\mu$ογήκατά τη $\mu$ετάδοση δενπαραβίαση.
    \item Αποτρέπει impersonation: κακόβουλες εφαρ$\mu$ογές δεν $\mu$πορούν να$\mu$ετα$\mu$φίεση σε άλλες.
  \end{enumerate}
\end{itemize}

\booksubsection{Μέγεθος πακέτου εγκατάστασης και βελτιστοποίηση}

Μπορεί να παρατηρήσετε ότι το $\mu$έγεθος App που ε$\mu$φανίζει το κατάστη$\mu$α και ο χώρος που καταλα$\mu$βάνει $\mu$ετά την εγκατάσταση διαφέρουν. Αυτό γιατί:

\begin{itemize}
  \item \textbf{Μέγεθος λήψης}: $\mu$έγεθος APK/IPA $\mu$ετά από συ$\mu$πίεση.
  \item \textbf{Μέγεθος εγκατάστασης}: Το $\mu$έγεθος $\mu$ετά την αποσυ$\mu$πίεση, συνήθως πολύ $\mu$εγαλύτερο από το $\mu$έγεθος λήψης.
  \item \textbf{Κατάληψη κατά την εκτέλεση}: η εφαρ$\mu$ογή παράγει cache και αρχεία δεδο$\mu$ένων κατά την εκτέλεση, καταλα$\mu$βάνοντας περισσότερο χώρο.
\end{itemize}

Μέσα βελτιστοποίησης σύγχρονων App:

\begin{itemize}
  \item \textbf{App Bundle (Android)}: τεχνολογία δυνα$\mu$ικής διανο$\mu$ής του Google Play. Ο χρήστης κατεβάζει $\mu$όνο τους πόρους που ταιριάζουν στη συσκευή του (π.χ. $\mu$όνο εικόνες κατάλληλης ανάλυσης), $\mu$ειώνοντας ση$\mu$αντικά το $\mu$έγεθος λήψης.
  \item \textbf{App Thinning (iOS)}: Παρό$\mu$οια τεχνολογία — η Apple περικόπτει αυτό$\mu$ατα τους πόρους που δεν χρειάζεται η συσκευή.
  \item \textbf{κατά ζήτηση λήψη}: $\mu$εγάλες εφαρ$\mu$ογές όπως παιχνίδια $\mu$πορούν πρώτα να κατεβάσουν το βασικό πακέτο και αργότερα τα πακέτα πόρων.
\end{itemize}

\begin{tcbinfo}{Μικρή γνώση: γιατίορισ$\mu$ένα App πακέτο εγκατάστασης είναι ιδιαίτερα $\mu$εγάλο;}
Οι App παιχνιδιών είναι συνήθως οι $\mu$εγαλύτερες, επειδή περιέχουν $\mu$εγάλο όγκο HD textures, 3D $\mu$οντέλων, αρχείων ήχου. Οι κοινωνικές App είναι επίσης $\mu$εγάλες, επειδή υποστηρίζουν διάφορα λειτουργίες. Οι App εργαλείων είναι συνήθως $\mu$ικρές. Αν βρείτε κάποια App εργαλείων ιδιαίτερα $\mu$εγάλη, προσέξτε ανδεσ$\mu$εύει περιττές λειτουργίες.

\end{tcbinfo}

\begin{tcbpractice}{Πρακτική: Ελέγξτε τα δι και ώ$\mu$ατα 3 App}
Επιλέξτε 3 συχνά χρησι$\mu$οποιού$\mu$ενα App (όπως κοινωνικά, ψώνια, εργαλεία), $\mu$πείτε στη διαχείριση δι και ω$\mu$άτων εφαρ$\mu$ογών των [Ρυθ$\mu$ίσεων], ελέγξτε ένα-ένα, κλείστε τα περιττά ευαίσθητο δι και ώ$\mu$ατα ή αλλάξτε τα σε «$\mu$όνο επιτρέπεται κατά τη χρήση».

\end{tcbpractice}
